\documentclass{chsmc}
\chsmcsetup{
  year = 2019,
  part = 1,
  type = B,
  show-answers = false,
}
\begin{document}

\section{填空题：本大题共 8 小题，每小题 8 分，共 64 分.}

\begin{enumerate}
  \item 已知实数集合 $\{1, 2, 3, x\}$ 的最大元素等于该集合的所有元素之和，
    则 $x$ 的值为 \blankline
  \item 若平面向量 $\vec{a} = (2^m, -1)$ 与
    $\vec{b} = (2^m - 1, 2^{m+1})$ 垂直，其中 $m$ 为实数，
    则 $\vec{a}$ 的模为 \blankline
  \item 设 $\alpha, \beta \in (0,\pi)$，$\cos\alpha$, $\cos\beta$
    是方程 $5x^2 - 3x - 1 = 0$ 的两根，则 $\sin\alpha\sin\beta$ 的值为 \blankline
  \item 设三棱锥 $P$-$ABC$ 满足 $PA = PB = 3$，$AB = BC = CA = 2$，
    则该三棱锥的体积的最大值为 \blankline
  \item 将 $5$ 个数 $2, 0, 1, 9, 2019$ 按任意次序排成一行，
    拼成一个 $8$ 位数（首位不为 $0$），则产生的不同的 $8$ 位数的个数为 \blankline
  \item 设整数 $n > 4$，$(x + 2\sqrt{y} - 1)^n$ 的展开式中 $x^{n-4}$ 与 $xy$
    两项的系数相等，则 $n$ 的值为 \blankline
  \item 在平面直角坐标系中，若以 $(r+1,0)$ 为圆心、$r$ 为半径的圆上存在一点
    $(a,b)$ 满足 $b^2 \geqslant 4a$，则 $r$ 的最小值为 \blankline
  \item 设等差数列 $\{a_n\}$ 的各项均为整数，首项 $a_1 = 2019$，
    且对任意正整数 $n$，总存在正整数 $m$，使得
    $a_1 + a_2 + \cdots + a_n = a_m$. 这样的数列 $\{a_n\}$ 的个数为 \blankline
\end{enumerate}

\section{解答题：本大题共 3 个小题，满分 56 分. 解答应写出文字说明、证明过程或演算步骤.}

\begin{enumerate}[resume]
  \item (本题满分 16 分) 在椭圆 $\Gamma$ 中，$F$ 为一个焦点，$A$, $B$ 为两个顶点.
    若 $|FA| = 3$，$|FB| = 2$，求 $|AB|$ 的所有可能值.
  \item (本题满分 20 分) 设 $a$, $b$, $c$ 均大于 $1$，满足
    $\begin{cases}
      \lg a + \log_b c = 3, \\
      \lg b + \log_a c = 4.
    \end{cases}$
    求 $\lg a \cdot \lg c$ 的最大值.
  \item (本题满分 20 分) 设复数数列 $\{z_n\}$ 满足：$|z_1| = 1$，
    且对任意正整数 $n$，均有 $4z_{n+1}^2 + 2z_n z_{n+1} + z_n^2 = 0$.
    证明：对任意正整数 $m$，均有
    $|z_1 + z_2 + \cdots + z_m| < \dfrac{2\sqrt{3}}{3}$.
\end{enumerate}

\end{document}